\chapter{Evaluation}

\section{Experimental configuration}

The frozen evaluation uses the repository's offline benchmark path with 256 synthetic scenes,
concurrency eight, the checked-in STAC fixture as a metadata template, temporary SQLite storage,
and temporary provenance storage. It exercises the actual validation, transformation,
bounded-ingestion, catalog, quality, and report code rather than a separate mock benchmark
implementation. No external provider, credential, charting package, or GIS dependency is
required. The resulting JSON and Markdown artifacts are checked in under
\texttt{docs/experiments/} and form the comparison point for later live runs.

\section{Frozen empirical baseline}

Table~\ref{tab:baseline} transcribes the publication-ready report in
\texttt{docs/experiments/experiment\_report.md}. Measurements are retained at their recorded
precision; elapsed time and latency describe the Windows development environment and should be
compared only with equivalent hardware, filesystem, Python version, and fixture.

\begin{table}[htbp]
\centering
\caption{Frozen 256-scene synthetic experiment baseline.}
\label{tab:baseline}
\begin{tabular}{lr}
\toprule
Metric & Value \\
\midrule
Scenes ingested & 256 \\
Elapsed seconds & 1.159317 \\
Throughput (items/sec) & 220.819653 \\
Failures & 0 \\
Retries & 0 \\
Payload bytes & 106130 \\
Catalog scenes & 256 \\
Usable scenes & 256 \\
Usable scene ratio & 1.000000 \\
Mean cloud cover & 4.200000 \\
Asset completeness ratio & 1.000000 \\
Spatial query latency (ms) & 9.117000 \\
Failure recovery ratio & 1.000000 \\
Payload megabytes & 0.106130 \\
\bottomrule
\end{tabular}
\end{table}

\section{Ingestion throughput and completeness}

The run ingested all 256 generated scenes in 1.159317 seconds, yielding 220.819653 items per
second. Throughput is calculated as
\[
\operatorname{throughput}=\frac{N_{\mathrm{ingested}}}{T_{\mathrm{elapsed}}}
=\frac{256}{1.159317}=220.819653\ \mathrm{items\,s^{-1}}.
\]
This is an end-to-end metadata-path measurement: validation, transformation, provenance writes,
SQLite indexing, and report accounting are included. It is therefore more informative for this
research question than a provider-only HTTP rate, but it is not a claim about raster download
throughput.

The catalog count equals the ingestion count (256), and failures and retries are both zero.
That equality indicates that the benchmark completed without a lost valid scene or a retry-induced
duplicate. The recovery ratio is also 1.000000, consistent with the fixture path having no
unrecoverable failures. These values demonstrate baseline correctness, not resilience under a
real provider outage; fault injection is evaluated separately by the test suite.

\section{Quality, payload, and storage footprint}

All 256 scenes are classified as usable, producing a usable-scene ratio of 1.000000. Mean cloud
cover is 4.200000 in the fixture metadata, while asset completeness is 1.000000. The quality
results show that the benchmark fixtures satisfy the intended schema and analytical prerequisites.
They should not be interpreted as an estimate of the cloud distribution of Sentinel-2 acquisitions
over Kvarken.

The serialized payload footprint is 106130 bytes, or 0.106130 megabytes under the report's
conversion. This is metadata-scale evidence: it measures the canonical scene payloads and not
the potentially multi-gigabyte raster assets referenced by their URLs. The small footprint
supports the educational design, in which a complete ingestion and query experiment can run
without downloading imagery.

\section{Spatial query latency}

The recorded spatial query latency is 9.117000 ms. The query path first applies indexed
bounding-box predicates and then performs exact regional checks, so the value reflects the
metadata catalog and geometry path rather than a scan of raster pixels. It is low enough for
interactive educational queries in the tested environment, while its single-run nature prevents
claims about tail latency or production concurrency.

The result also illustrates why the architecture reports query latency separately from ingestion
throughput. Ingestion measures a pipeline over 256 records; spatial latency measures a query
operation over the resulting catalog. Conflating the two would hide whether an observed slowdown
comes from provider processing, persistence, or regional filtering.

\section{Validation and recovery evidence}

The release baseline passes 53 offline tests. The suite covers typed schema rejection, pagination,
timeouts, HTTP 429 recovery, OAuth2 expiry and invalid token payloads, provenance integrity,
bounded worker limits, quality filters, raster ranges and NDVI, HTTP queries, artifact retention,
multi-threaded SQLite contention, WAL durability after reopening, and deterministic reports.
The worker stress tests observe the maximum active task count rather than merely checking the
final row count. The catalog tests accept only committed intermediate snapshots during contention
and verify that committed state survives a reopen.

Ruff lint and format checks pass, and \texttt{python -m kvarken\_eo verify-health} completes the
fixture-to-provenance-to-catalog maintenance path. Together, these checks provide evidence for
implementation integrity, while the benchmark provides an empirical reference. Neither replaces
evaluation against authenticated CDSE data or representative raster windows.

\section{Threats to validity}

\subsection{Internal validity}

The benchmark uses synthetic scenes derived from a checked-in fixture, so the measured zero
failure rate is partly a consequence of controlled inputs. Timing can also be affected by Windows
filesystem cache state, interpreter scheduling, temporary-directory placement, and background
processes. The mitigation is reproducibility: the fixture, report, concurrency setting, and
runtime path are versioned, and the 53-test suite injects timeouts, 429 responses, malformed token
payloads, concurrent readers, and bounded-worker probes. Repeated runs should be interpreted as
comparisons under controlled conditions, not as independent population samples.

\subsection{External validity}

The results do not establish behavior for all EO missions, CDSE collections, network conditions,
or raster sizes. Sentinel-2-like metadata is representative for the architecture but does not
capture provider-specific schema extensions, authentication policy changes, regional datacubes, or
large binary transfers. The live runner and dated artifacts provide the intended extension point:
future authenticated observations can be added without overwriting the frozen baseline. Full
GeoTIFF decoding, distributed storage, and multi-process catalogs remain outside this evaluation.

\subsection{Construct validity}

Throughput, payload bytes, usable-scene ratio, and spatial latency are operational proxies, not
complete measures of scientific usefulness, sustainability, or user learning. A successful
metadata ingestion does not prove pixel-level radiometric correctness. Likewise, a 9.117 ms query
does not measure end-user latency through a deployed network service. The study addresses these
construct limits by stating exactly what each metric includes and by keeping raster transport,
NDVI behavior, quality rules, and HTTP API checks as separate testable slices.

\subsection{Harness reliability}

The test harness is itself a controlled artifact. Injectable transports make failure modes
deterministic; temporary directories isolate provenance and SQLite state; worker-bound tests
measure concurrency directly; and WAL tests verify both valid snapshots during contention and
durability after reopening. The health command exercises the principal offline path in one
operation. These controls reduce accidental green results, but they cannot remove all timing
variance or substitute for long-running production observation.

\section{Interpretation}

The evidence supports the thesis claim that a dependency-free, typed, bounded, provenance-aware
metadata pipeline can provide a reproducible regional EO baseline suitable for research and
education. The result is deliberately narrower than a production-scale data platform: its value
lies in explicit contracts and an auditable starting point. Live CDSE measurements, representative
raster analyses, and comparative deployments are the appropriate next experiments, with the
256-scene artifact retained as a stable regression reference.
